\subsection{Why independently marked disks do not globalize the local sources}

\begin{proposition}[D022: \status{Proved}; integral elementary-source obstruction]
\label{sa08:D022}
Let $S=\Gamma_{\Z}(T)$. For an arbitrary set $E$, choose
$d_\alpha>0$ and $a_\alpha\in\Z$ for every $\alpha\in E$, and put
\[
 R=\Gamma_{\Z}(U_\alpha:\alpha\in E),\qquad
 |U_\alpha|=d_\alpha,\qquad
 U_\alpha\longmapsto a_\alpha T^{(d_\alpha)}.
\]
Let $\Phi_\alpha$ be the standard contractible simplicial marked
module with its generators assigned weight $d_\alpha$. Define
\[
 G=S\otimes_R
 \mathrm{DR}\left(\bigoplus_{\alpha\in E}\Phi_\alpha;
                         \{U_\alpha\}_{\alpha\in E}\right),
\]
where the differential kills each marked line and is the usual
PD Poincare differential on a free complement. Then $G$ is an
eventually admissible theory over $\Z$ if and only if $E$ consists
of one element, $d_\alpha=1$, and $a_\alpha=1$ or $-1$.
Every other choice fails eventual horizontal acyclicity already
in exterior degree zero.
\end{proposition}
\begin{proof}
The marked lines split in each simplicial degree. Thus the displayed
de Rham complex is $R$ tensored with a free PD Poincare complex;
after base change it resolves $S$ vertically. Its terms are free
abelian groups, possibly of infinite rank. Furthermore,
$\Gamma(\bigoplus_\alpha\Phi_\alpha)$ is a simplicial resolution
of $\Z=R/R^+$ by $R$-free modules. Indeed, the direct sum of the
marked disks is a contractible simplicial module, and positive
divided powers preserve its simplicial contraction. Hence
\[
 \pi_i(G^0)=\operatorname{Tor}_i^R(S,\Z).
\]
Suppose $G^0$ is acyclic in all weights $w\ge W$.
Each such normalized weight complex is a nonnegative acyclic
complex of projective abelian groups. Its cycle sequences split
inductively from degree zero, giving a chain contraction.
Reduction modulo any prime $p$ therefore preserves acyclicity in
these weights.

Write $R_p=R\otimes\mathbb F_p$ and $S_p=S\otimes\mathbb F_p$.
Vanishing of $\pi_0$ in high weights says that
$S_p/f(R_p^+)S_p$ is bounded in weight. The indecomposables of
$S_p$ have basis $T^{(p^s)}$, $s\ge0$; those of $R_p$ have basis
$\gamma_{p^r}(U_\alpha)$. For large $s$, an expression for
$T^{(p^s)}$ in the ideal $f(R_p^+)S_p$ must contain a summand whose
image in $S_p^+/(S_p^+)^2$ is nonzero. Its multiplier has weight
zero, and its divided-power generator has weight
$d_\alpha p^r=p^s$. Necessarily some $\alpha$ has
\[
 d_\alpha=N=p^e,\qquad p\nmid a_\alpha.
\]
For this $\alpha$ the algebra $S_p$ is free over
$R_\alpha=\Gamma_{\mathbb F_p}(U_\alpha)$, with basis
$T^{(r)}$, $0\le r<N$. To verify this directly, the coefficient of
$T^{(mN+r)}$ in $U_\alpha^{(m)}T^{(r)}$ is
\[
 a_\alpha^m\frac{(mN)!}{m!(N!)^m}\binom{mN+r}{r}.
\]
Legendre's valuation formula gives valuation zero for the factorial
quotient, and the base-$p$ digits of $mN$ and $r$ do not overlap,
so the binomial coefficient is also a unit modulo $p$. Thus
\[
 S_p\otimes_{R_\alpha}\mathbb F_p=S_p/I_N,
 \qquad I_N=\bigoplus_{w\ge N}S_p^w.
\]
Decomposing $R_p=R_\alpha\otimes R_{\mathrm{rest}}$, freeness and
associativity of derived tensor give
\[
 S_p\otimes_{R_p}^{\mathbf L}\mathbb F_p
 \simeq (S_p/I_N)\otimes_{R_{\mathrm{rest}}}^{\mathbf L}\mathbb F_p.
\]

Suppose there is another index $\beta$. The divided-power basis
gives the algebra decomposition
\[
 \Gamma_{\mathbb F_p}(U_\beta)
   \cong\bigotimes_{r\ge0}
       \mathbb F_p[y_{\beta,r}]/(y_{\beta,r}^{p}),
 \qquad |y_{\beta,r}|=d_\beta p^r.
\]
The generator $y_{\beta,r}$ is a nonzero scalar multiple of
$\gamma_{p^r}(U_\beta)$; the required scalars are units by the
base-$p$ digit factorial formula. Split all these tensor factors
of $R_{\mathrm{rest}}$ into
$R_{\mathrm{rest}}=L\otimes H$, where $L$ consists of the factors
whose generators have weight less than $N$ and $H$ of those whose
generators have weight at least $N$. Every generator of $H$ acts
as zero on $V=S_p/I_N$. Tensoring free resolutions over the field
$\mathbb F_p$ gives
\[
 \operatorname{Tor}_*^{R_{\mathrm{rest}}}(V,\mathbb F_p)
 \cong
 \operatorname{Tor}_*^L(V,\mathbb F_p)\otimes_{\mathbb F_p}
 \operatorname{Tor}_*^H(\mathbb F_p,\mathbb F_p).
\]
The weight-zero part of the first factor in homological degree
zero is $\mathbb F_p$: the augmentation
$V/L^+V\to\mathbb F_p$ is surjective. The second factor has
\[
 \operatorname{Tor}_1^H(\mathbb F_p,\mathbb F_p)=H^+/(H^+)^2,
\]
which contains the independent classes of $y_{\beta,r}$ for all
sufficiently large $r$. Their weights are unbounded. We obtain
nonzero $\pi_1(G_p^0)$ in unbounded weights, a contradiction.
This also covers infinitely many marked generators: infinite
tensor products and the resolutions involved are filtered
colimits, and the displayed classes are already present in finite
subsets of the factors.

The index set therefore has at most one element. It cannot be
empty, since then $G^0=S$ is constant and its positive weights
have nonzero $\pi_0$. For the unique element, the earlier argument
holds at every prime: $d$ must be a power of every prime, hence
$d=1$, and $a$ must be divisible by no prime, hence $a=\pm1$.
Conversely these data give the standard theory, after a change of
sign of its marked generator if necessary. It is admissible.
\end{proof}

\begin{corollary}[D023: \status{Proved}; a necessary departure from elementary sources]
\label{sa08:D023}
Let $N>1$ and let $A_N\subset F$ be the integral standard tail
subtheory: its weights below $N$ are just the coefficient algebra
$S$, and its weights at least $N$ agree with the standard theory
$F$. No admissible source of the form in
Statement~\ref{sa08:D022} maps to $A_N$.
\end{corollary}
\begin{proof}
The positive weights add under multiplication and multiply under
divided powers, so $A_N$ is a simplicial PD sub-DGA of $F$. It
resolves $S$ in every weight and is flasque in exterior degree $q$
when $w\ge\max(N,q+1)$, hence is admissible. Its weight-one,
exterior-degree-zero simplicial module is the constant $\Z T$.
Statement~\ref{sa08:D022} leaves only the standard marked disk of
weight one as a possible elementary source. A map to $A_N$ would
send its marked vector to $T$ or $-T$. But the marked vector is a
normalized simplicial boundary in that disk, whereas a constant
simplicial module has zero normalized differential. No such map
exists.
\end{proof}

The supplementary program
\texttt{code/D\_elementary\_source\_checks.py} checks the free-basis
coefficients for $p=2,3,5$, $N=p^e$ with $0\le e\le2$, and
$0\le m\le8$. It also checks the reduced bar differential for one
additional zero-acting truncated-polynomial factor. Its saved output
records $477$ assertions. The unbounded persistence needed in the
proof is proved above, not inferred from these finite checks.

These statements do not refute common-source uniqueness in the
full category: $A_N$ itself maps to $F$ and is a non-elementary
source for that pair. They prove that adding arbitrarily many
independently marked disks cannot globalize the $p$-local argument.
An integral proof must construct a source with additional
relations or use a different comparison mechanism.
